\documentclass[12pt]{article}
\usepackage[utf8]{inputenc}
\usepackage[T1]{fontenc}
\usepackage{hyperref}
\title{Foia Foia Language Resources for K-12 Curriculum Development}
\author{Compiled by Claude Sonnet 5.5 (claude-sonnet-5-5) for Andr\'as Kornai}
\date{}
\begin{document}
\maketitle
\section{Introduction}
Foia Foia (ISO 639-3 ffi; Glottolog foia1236) is a Papuan language of Papua New Guinea, spoken near the Omati River mouth in Gulf Province, with villages also reported in Western Province. Wikipedia gives about 180 speakers (2000 census); Wikitongues rates it threatened. Resources are very thin.
\section{Primary Educational Resources}
None found.
\section{Linguistic and Reference Materials}
\begin{itemize}
\item Wikipedia, Foia Foia language: \url{https://en.wikipedia.org/wiki/Foia_Foia_language}
\item Wikitongues language page: \url{https://wikitongues.org/languages/ffi/}
\end{itemize}
\section{Religious and Community Resources}
\begin{itemize}
\item Joshua Project language page (translation listed only as started): \url{https://joshuaproject.net/languages/ffi}
\item Joshua Project people group page (its location details conflict with Wikipedia and should be treated with caution): \url{https://joshuaproject.net/people_groups/19682/PP}
\end{itemize}
\section{Digital and Audio Resources}
Global Recordings Network reports no recordings: \url{https://globalrecordings.net/en/language/ffi}
\section{Curriculum Development Context}
No school use or literacy materials were found. Any curriculum work would begin with documentation.
\end{document}